\documentclass[12pt]{article}
\usepackage[utf8]{inputenc}
\usepackage[russian]{babel}
\usepackage{amsmath,amssymb}
\usepackage{graphics}
\usepackage{pbox}
\usepackage[x11names]{xcolor}
\definecolor{brightmaroon}{rgb}{0.76, 0.13, 0.28}
\definecolor{royalazure}{rgb}{0.0, 0.22, 0.66}
\usepackage[colorlinks=true,linkcolor=royalazure]{hyperref}
\usepackage{tikz, tkz-fct, pgfplots}
\usetikzlibrary{arrows}
\usepackage{geometry}
\geometry{
	a4paper,
	total={170mm,257mm},
	left=20mm,
	top=20mm
} 
\usepackage[labelsep=period]{caption}

% ----------------- Commands ----------------- 
\newcommand{\eps}{\varepsilon}
\newcommand\tline[2]{$\underset{\text{#1}}{\text{\underline{\hspace{#2}}}}$}

% ----------------- Set graphics path ----------------- 
\graphicspath{{img/}}

\begin{document}
\pagestyle{empty}
\centerline{\large Министерство науки и высшего образования}	
\centerline{\large Федеральное государственное бюджетное образовательное}
\centerline{\large учреждение высшего образования}
\centerline{\large ``Московский государственный технический университет}
\centerline{\large имени Н.Э. Баумана}
\centerline{\large (национальный исследовательский университет)''}
\centerline{\large (МГТУ им. Н.Э. Баумана)}
\hrule

\vspace{0.5cm}

\begin{figure}[h]
\center
\includegraphics[height=0.35\linewidth]{bmstu-logo-small.png}
\end{figure}

\begin{center}
	\large	
	\begin{tabular}{c}
		Факультет ``Фундаментальные науки'' \\
		Кафедра ``Высшая математика''		
	\end{tabular}
\end{center}

\vspace{0.5cm}

\begin{center}
	\LARGE \bf	
	\begin{tabular}{c}
		\textsc{Отчёт} \\
		по учебной практике \\
		за 1 семестр 2020---2021 гг.
	\end{tabular}
\end{center}

\vspace{0.5cm}

\begin{center}
	\large
	\begin{tabular}{p{5.3cm}ll}
		\pbox{5.45cm}{
			Руководитель практики,\\
			ст. преп. кафедры ФН1} 	& \tline{\it(подпись)}{5cm} & Кравченко О.В. \\[0.5cm]
		студент группы ФН1--11 		& \tline{\it(подпись)}{5cm} & Ф.И.О.
	\end{tabular}
\end{center}

\vfill
\begin{center}
	\large	
	\begin{tabular}{c}
		Москва, \\
		2020 г.
	\end{tabular}
\end{center}

\newpage	
\tableofcontents

\newpage
\section{Цели и задачи практики}	
\subsection{Цели}
--- развитие компетенций, способствующих успешному освоению материала бакалавриата и необходимых в будущей профессиональной деятельности.

\subsection{Задачи}
\begin{enumerate}
\item Знакомство с программными средствами, необходимыми в будущей профессиональной деятельности.
\item Развитие умения поиска необходимой информации в специальной литературе и других источниках.
\item Развитие навыков составления отчётов и презентации результатов.
\end{enumerate}

\subsection{Индивидуальное задание}	
\begin{enumerate}
\item Изучить способы отображения математической информации в системе вёртски \LaTeX.
\item Изучить возможности  системы контроля версий \textsf{Git}.
\item Научиться верстать математические тексты, содержащие формулы и графики в системе \LaTeX.
Для этого, выполнить установку свободно распространяемого дистрибутива \textsf{TeXLive} и оболочки \textsf{TeXStudio}.
\item Оформить в системе \LaTeX типовые расчёты по курсе математического анализа согласно своему варианту.
\item Создать аккаунт на онлайн ресурсе \textsf{GitHub} и загрузить исходные \textsf{tex}--файлы 
и результат компиляции в формате \textsf{pdf}.
\end{enumerate} 

\newpage
\section{Отчёт}
Актуальность темы продиктована необходимостью владеть системой вёрстки \LaTeX и средой вёрстки \textsf{TeXStudio} для
отображения текста, формул и графиков. Полученные в ходе практики навыки могут быть применены при написании
курсовых проектов и дипломной работы, а также в дальнейшей профессиональной деятельности.

Ситема вёрстки \LaTeX содержит большое количество инструментов (пакетов), упрощающих отображение информации в различных 
сферах инженерной и научной деятельности. 

\newpage
\section{Индивидуальное задание}
\subsection{Элементарные функции и их графики.}
% ---------------------------- Problem 1----------------------------------
\subsubsection*{\center Задача № 1.}
{\bf Условие.~}
Найти область определения функции
$$
y(x) = \lg\biggl(\dfrac{x+4}{1-2x}\biggr).
$$	
{\bf Решение.~}	
$D_f: \dfrac{x+4}{1-2x}>0 \Leftrightarrow x \in (-4;\,-1/2).$

% ---------------------------- Problem 2----------------------------------
\subsubsection*{\center Задача № 2.}
{\bf Условие.~}
Исследовать функцию на чётность и нечётность
$$
y(x) = \ctg(\cos(\tg(x))).
$$	
{\bf Решение.~}	
$$
y(-x) = \ctg(\cos(\tg(-x))) = \ctg(\cos(-\tg(x))) = \ctg(\cos(\tg(x))) = y(x).
$$	
Отсюда, $y(x)$ --- чётная функция.


% ---------------------------- Problem 3----------------------------------
\subsubsection*{\center Задача № 3.}
{\bf Условие.~}
Используя элементарные преобразования, построить эскизы графиков функций (а)--(д).
$$
\begin{array}{cc}
\text{3(а):} & y(x) = -1 - \sin(2x + \pi/4), \\
\text{3(б):} & y(x) = |2 \sqrt[3]{x+5} - 1|, \\
\text{3(в):} & y(x) = 1 - \log_3|x+1|, \\
\text{3(г):} & y(x) = \dfrac{1}{3} 2^{2x+1} - \dfrac{4}{3}, \\[8pt]
\text{3(д):} & y(x) = \dfrac{3\pi}{8} + \dfrac{3}{2}\arctg(2x-3).
\end{array}
$$
{\bf Решение.~}
% ---------------- Problem 3а----------------
Последовательность элементарных преобразований графика функции 3(а).
\begin{center}
	\begin{tikzpicture}
	\begin{axis}[width=12cm, height=4cm, samples=250, domain=-2*pi:2*pi,legend pos=outer north east]
	\addplot[blue, dotted, ultra thick, samples=250, domain=-2*pi:2*pi]{sin(deg(x))};
	\addlegendentry{$\sin(x)$}
	\addplot[blue, ultra thick, samples=250, domain=-2*pi:2*pi]{sin(deg(x+pi/8))};
	\addlegendentry{$\sin(x+\pi/8)$}
	\end{axis}
	\end{tikzpicture}
	\begin{tikzpicture}
	\begin{axis}[width=12cm, height=4cm, samples=250, domain=-2*pi:2*pi,legend pos=outer north east]
	\addplot[blue, dotted, ultra thick, samples=250, domain=-2*pi:2*pi]{sin(deg(x+pi/8))};
	\addlegendentry{$\sin(x+\pi/8)$}
	\addplot[blue, ultra thick, samples=250, domain=-2*pi:2*pi]{sin(2*deg(x+pi/8))};
	\addlegendentry{$\sin(2x+\pi/4)$}
	\end{axis}
	\end{tikzpicture}
	\begin{tikzpicture}
	\begin{axis}[width=12cm, height=4cm, samples=250, domain=-2*pi:2*pi,legend pos=outer north east]
	\addplot[blue, dotted, ultra thick, samples=250, domain=-2*pi:2*pi]{sin(2*deg(x+pi/8))};
	\addlegendentry{$\sin(2x+\pi/4)$}
	\addplot[blue, ultra thick, samples=250, domain=-2*pi:2*pi]{-sin(2*deg(x+pi/8))};
	\addlegendentry{-$\sin(2x+\pi/4)$}
	\end{axis}
	\end{tikzpicture}
	\begin{tikzpicture}
	\begin{axis}[width=12cm, height=4cm, samples=250, domain=-2*pi:2*pi,legend pos=outer north east]
	\addplot[blue, dotted, ultra thick, samples=250, domain=-2*pi:2*pi]{-sin(2*deg(x+pi/8))};
	\addlegendentry{$-\sin(2x+\pi/4)$}
	\addplot[blue, ultra thick, samples=250, domain=-2*pi:2*pi]{-1-sin(2*deg(x+pi/8))};
	\addlegendentry{$-1-\sin(2x+\pi/4)$}
	\end{axis}
	\end{tikzpicture}
\end{center}

% ---------------- Problem 3б----------------
Последовательность элементарных преобразований графика функции 3(б).
\begin{center}
	\begin{tikzpicture}
	\begin{axis}[width=12cm, height=6cm, samples=250, domain=-6:6,legend pos=outer north east]
	\addplot[blue, dotted, ultra thick]{sign(x)*abs(x)^(1/3)};
	\addlegendentry{$\sqrt[3]{x}$}
	\addplot[blue, ultra thick]{sign(x+5)*abs(x+5)^(1/3)};
	\addlegendentry{$\sqrt[3]{x+5}$}	
	\end{axis}
	\end{tikzpicture}
	\begin{tikzpicture}
	\begin{axis}[width=12cm, height=6cm, samples=250, domain=-6:6,legend pos=outer north east]
	\addplot[blue, dotted, ultra thick]{sign(x+5)*abs(x+5)^(1/3)};
	\addlegendentry{$\sqrt[3]{x+5}$}	
	\addplot[blue, ultra thick]{sign(x+5)*abs(x+5)^(1/3)-1/2};
	\addlegendentry{$\sqrt[3]{x+5}-\frac12$}
	\end{axis}
	\end{tikzpicture}
	\begin{tikzpicture}
	\begin{axis}[width=12cm, height=6cm, samples=250, domain=-6:6,legend pos=outer north east]
	\addplot[blue, dotted, ultra thick]{sign(x+5)*abs(x+5)^(1/3)-1/2};
	\addlegendentry{$\sqrt[3]{x+5}-\frac12$}
	\addplot[blue, ultra thick]{2*(sign(x+5)*abs(x+5)^(1/3)-1/2)};
	\addlegendentry{$2\sqrt[3]{x+5}-1$}	
	\end{axis}
	\end{tikzpicture}
	\begin{tikzpicture}
	\begin{axis}[width=12cm, height=6cm, samples=250, domain=-6:6,legend pos=outer north east]
	\addplot[blue, dotted, ultra thick]{2*sign(x+5)*abs(x+5)^(1/3)-1};
	\addlegendentry{$2\sqrt[3]{x+5}-1$}
	\addplot[blue, ultra thick]{abs(2*(sign(x+5)*abs(x+5)^(1/3)-1/2))};
	\addlegendentry{$|2\sqrt[3]{x+5}-1|$}	
	\end{axis}
	\end{tikzpicture}
\end{center}

% ---------------- Problem 3в----------------
Последовательность элементарных преобразований графика функции 3(в).
\begin{center}
	\begin{tikzpicture}
	\begin{axis}[width=12cm, height=6cm, samples=250, domain=-6:6,legend pos=outer north east]
	\addplot[blue, ultra thick]{ln(abs(x))/ln(3)};
	\addlegendentry{$\log_3{|x|}$}	
	\addplot[red, dotted, ultra thick]{ln(x)/ln(3)};
	\addlegendentry{$\log_3{x}$}	
	\end{axis}
	\end{tikzpicture}
	\begin{tikzpicture}
	\begin{axis}[width=12cm, height=6cm, samples=250, domain=-6:6,legend pos=outer north east]
	\addplot[blue, dotted, ultra thick]{ln(abs(x))/ln(3)};
	\addlegendentry{$\log_3{|x|}$}	
	\addplot[blue, ultra thick]{ln(abs(x+1))/ln(3)};
	\addlegendentry{$\log_3{|x+1|}$}	
	\end{axis}
	\end{tikzpicture}
	\begin{tikzpicture}
	\begin{axis}[width=12cm, height=6cm, samples=250, domain=-6:6,legend pos=outer north east]
	\addplot[blue, dotted, ultra thick]{ln(abs(x+1))/ln(3)};
	\addlegendentry{$\log_3{|x+1|}$}	
	\addplot[blue, ultra thick]{-ln(abs(x+1))/ln(3)};
	\addlegendentry{$-\log_3{|x+1|}$}	
	\end{axis}
	\end{tikzpicture}
	\begin{tikzpicture}
	\begin{axis}[width=12cm, height=6cm, samples=250, domain=-6:6,legend pos=outer north east]
	\addplot[blue, dotted, ultra thick]{-ln(abs(x+1))/ln(3)};
	\addlegendentry{$-\log_3{|x+1|}$}	
	\addplot[blue, ultra thick]{1-ln(abs(x+1))/ln(3)};
	\addlegendentry{$1-\log_3{|x+1|}$}	
	\end{axis}
	\end{tikzpicture}
\end{center}

% ---------------- Problem 3г----------------
Последовательность элементарных преобразований графика функции 3(г).
\begin{center}
	\begin{tikzpicture}
	\begin{axis}[width=12cm, height=6cm, samples=250, domain=-3:3,legend pos=outer north east]
	\addplot[blue, dotted, ultra thick]{2^x};
	\addlegendentry{$2^x$}	
	\addplot[blue, ultra thick]{2^(x+1/2)};
	\addlegendentry{$2^{x+1/2}$}	
	\end{axis}
	\end{tikzpicture}
	\begin{tikzpicture}
	\begin{axis}[width=12cm, height=6cm, samples=250, domain=-3:3,legend pos=outer north east]
	\addplot[blue,dotted, ultra thick]{2^(x+1/2)};
	\addlegendentry{$2^{(x+1/2)}$}	
	\addplot[blue, ultra thick]{2^(2*x+1)};
	\addlegendentry{$2^{2x+1}$}	
	\end{axis}
	\end{tikzpicture}
	\begin{tikzpicture}
	\begin{axis}[width=12cm, height=6cm, samples=250, domain=-3:3,legend pos=outer north east]
	\addplot[blue, dotted, ultra thick]{2^(2*x+1)};
	\addlegendentry{$2^{2x+1}$}	
	\addplot[blue, ultra thick]{(1/3)*2^(2*x+1)};
	\addlegendentry{$\frac{1}{3}2^{2x+1}$}	
	\end{axis}
	\end{tikzpicture}
	\begin{tikzpicture}
	\begin{axis}[width=12cm, height=6cm, samples=250, domain=-3:3,legend pos=outer north east]
	\addplot[blue, dotted, ultra thick]{(1/3)*2^(2*x+1)};
	\addlegendentry{$\frac{1}{3}2^{2x+1}$}	
	\addplot[blue, ultra thick]{(1/3)*2^(2*x+1)-(4/3)};
	\addlegendentry{$\frac{1}{3}2^{2x+1}-\frac43$}	
	\end{axis}
	\end{tikzpicture}
\end{center}

% ---------------- Problem 3д----------------
Последовательность элементарных преобразований графика функции 3(д).
\begin{center}
	\begin{tikzpicture}
	\begin{axis}[width=12cm, height=6cm, samples=250, domain=-8:8,legend pos=outer north east]
	\addplot[blue, dotted, ultra thick]{rad(atan(x))};
	\addlegendentry{$\arctg{x}$}	
	\addplot[blue, ultra thick]{rad(atan(x-3/2))};
	\addlegendentry{$\arctg(x-3/2)$}	
	\end{axis}
	\end{tikzpicture}
	\begin{tikzpicture}
	\begin{axis}[width=12cm, height=6cm, samples=250, domain=-8:8,legend pos=outer north east]
	\addplot[blue, dotted, ultra thick]{rad(atan(x-3/2))};
	\addlegendentry{$\arctg{(x-3/2)}$}	
	\addplot[blue, ultra thick]{rad(atan(2*x-3))};
	\addlegendentry{$\arctg(2x-3)$}	
	\end{axis}
	\end{tikzpicture}
	\begin{tikzpicture}
	\begin{axis}[width=12cm, height=6cm, samples=250, domain=-8:8,legend pos=outer north east]
	\addplot[blue, dotted, ultra thick]{rad(atan(2*x-3))};
	\addlegendentry{$\arctg{2x-3}$}	
	\addplot[blue, ultra thick]{1.5*rad(atan(2*x-3))};
	\addlegendentry{$\frac32\arctg(2x-3)$}	
	\end{axis}
	\end{tikzpicture}
	\begin{tikzpicture}
	\begin{axis}[width=12cm, height=6cm, samples=250, domain=-8:8,legend pos=outer north east]
	\addplot[blue, dotted, ultra thick]{1.5*rad(atan(2*x-3))};
	\addlegendentry{$\frac32\arctg{2x-3}$}	
	\addplot[blue, ultra thick]{(3*pi/8) + 1.5*rad(atan(2*x-3))};
	\addlegendentry{$\frac{3\pi}{8} + \frac32\arctg(2x-3)$}	
	\end{axis}
	\end{tikzpicture}
\end{center}

% ---------------------------- Problem 4----------------------------------
\subsubsection*{\center Задача № 4.}
{\bf Условие.~}
Построить эскиз графика рациональной функции, найдя его асимптоты и исследуя 
расположение графика относительно оси абсцисс и асимптот (не используя предела)
$$
y(x) = \dfrac{2x^2-x-3}{x^2-4x+4}.
$$
{\bf Решение.~}	
Выделим целую часть
$$
y(x) = \dfrac{2x^2-x-3}{x^2-4x+4} = \dfrac{2x^2-8x+8+7x-11}{(x-2)^2} = \dfrac{2(x-2)^2+7x-11}{x^2-4x+4} =
2 + \dfrac{7x-11}{(x-2)^2} .
$$
Отсюда, $y = 2$ --- горизонтальная ассимптота, а также, $x = 2$ --- вертикальная ассимптота.
\begin{center}
	\begin{tikzpicture}
	\begin{axis}[
	width=12cm, height=12cm,
	axis lines = center,
	xlabel = $x$,
	ylabel = {$y$},
	xmax = {8.5},
	xmin = {-5.5},
	ymax = {10.5},
	ymin = {-4.5},
	restrict y to domain = -5:10,
	legend pos = outer north east
	]
	\addplot [
	domain=-5:8,
	samples=500,
	color=black,
	ultra thick
	]
	{(2*x^2-x-3)/(x^2-4*x+4)};
	%\addlegendentry{2 turning points}
	
	% Oblique asymptote at y=x
	\addplot[dashed] {2};
	% Vertical asymptote at x=2
	\draw[dashed] ({axis cs:2,0}|-{rel axis cs:0,0}) -- ({axis cs:2,0}|-{rel axis cs:0,1});
	\end{axis}
	\end{tikzpicture}
\end{center}

% ---------------------------- Problem 5----------------------------------
\subsubsection*{\center Задача № 5.}
{\bf Условие.~}
Построить эскиз графика сложной функции, используя различные элементарные приёмы
$$
y(x) = \biggl(\dfrac{3}{2}\biggr)^x\cdot\cos\biggl(\dfrac{\pi x}{2}\biggr).
$$
{\bf Решение.~}	График функции $y(x)=\cos(\pi x/2)$ получается элементарными преобразованиями из
графика функции $y(x)=\cos(x)$. В то же время, график функции $y(x)=(3/2)^x\cos(\pi x/2)$ получается
из графика функции $y(x)=\cos(\pi x/2)$  пропорциональным изменением ординаты каждой точки графика в $(3/2)^x$ раз для
каждого значения абсциссы $x$.
\begin{center}
	\begin{tikzpicture}
	\begin{axis}[width=12cm, height=6cm, samples=250, domain=-3.5:4,legend pos=outer north east]
	\addplot[blue, dotted, ultra thick, samples=250, domain=-3.5:4]{cos(deg(0.5*pi*x))};
	\addlegendentry{$\cos(\frac{\pi x}{2}$}
	\addplot[blue, ultra thick, samples=250, domain=-3.5:4]{((1.5)^x)*cos(deg(0.5*pi*x))};
	\addlegendentry{$(\frac32)^x\cos(\frac{\pi x}{2})$}
	\end{axis}
	\end{tikzpicture}
\end{center}

%=================================================================================================================================
\subsection{Пределы и непрерывность.}
% ---------------------------- Problem 1----------------------------------
\subsubsection*{\center Задача № 1.}
{\bf Условие.~}
Разложить в ряд Фурье заданную функцию $f(x)$, построить графики $f(x)$ и суммы ее ряда Фурье. Если не указывается, какой вид разложения в ряд необходимо представить, то требуетчя разложить функцию либо в общий тригонометрический ряд Фурье, либо следует выбрать оптимальный вид разложения в зависимости от данной функции.

\[
f(x)=\begin{cases}
-1, 	& 0 \leqslant x \leqslant 2,\\
x-3, 	& 2 \leqslant x \leqslant 3.
\end{cases}
\]
{\bf Решение.~}	
%График
\begin{center}
	\begin{tikzpicture}
	\begin{axis}[xmin=-1,	xmax=3.5, 	ymin=-1,	ymax=0.5,
	width=0.5\textwidth,
	height=0.4\textwidth,
	axis x line=middle,
	axis y line=middle, 
	every axis x label/.style={at={(current axis.right of origin)},anchor=west},
	every inner x axis line/.append style={|-latex'},
	every inner y axis line/.append style={|-latex'},
	minor tick num=1,			
	axis equal=true,
	xlabel=$x$, 
	ylabel=$y$,          
	samples=100,
	clip=true,
	]
	\addplot[color=black, line width=1.5pt,domain=0:2] {-1};
	\addplot[color=black, line width=1.5pt,domain=2:3]{\x-3};
	\addplot[thick,dashed] coordinates {(2,0) (2,-1)};
	\addplot[
	mark=*,
	mark options={fill=black, draw=black},
	only marks,
	] coordinates {(2, -1)};
	\end{axis}
	\end{tikzpicture}
\end{center}
\noindent
Построим общий тригонометрический ряд Фурье вида
$$
f(x)=\frac{a_0}{2}+\sum_{n=1}^\infty 
	\left(a_n\cos{(n\omega x)}+b_n\sin{(n\omega x)}\right),\quad\text{где}\,\omega=\frac{2\pi}{T},\,T=3.
$$
\noindent
Вычислим коэффициенты
$$
\begin{array}{rcl}
a_0 &=& \displaystyle\frac{2}{3}\left.\left(
\int\limits_0^2
(-1)\,dx + \int\limits_2^3
(x-3)\,dx \right) = 
\frac{2}{3}\left(
-2 + \left( \frac{x^2}{2}-3x\right)
\right|_2^3 \right) = -\frac{5}{3},												\\[12pt]
a_n &=& \displaystyle\frac{2}{3}\left(
	-\int\limits_0^2
	\cos \frac{2\pi nx}{3}\,dx + \int\limits_2^3
	(x-3)\cos \frac{2\pi nx}{3}\,dx \right) ={}									\\[12pt]
	&=& \displaystyle\frac{2}{3}\left(
	-\left.\frac{3}{2\pi n}\sin\frac{2\pi nx}{3} \right|_0^2
	+\left.\frac{3}{2\pi n}(x-3)\sin\frac{2\pi nx}{3} \right|_2^3 
	+\left.\frac{9}{(2\pi n)^2}(x-3)\cos\frac{2\pi nx}{3} \right|_2^3\right) = 	\\[12pt]
	&=& \displaystyle-\frac{3}{2\pi^2 n^2}\left(1-\cos\frac{4\pi n}{3} \right),	\\[12pt]
b_n &=& \displaystyle\frac{2}{3}\left(
	-\int\limits_0^2
	\sin \frac{2\pi nx}{3}\,dx + \int\limits_2^3
	(x-3)\sin \frac{2\pi nx}{3}\,dx \right) ={}									\\[12pt]
	&=& \displaystyle\frac{2}{3}\left(
	\left.\frac{3}{2\pi n}\cos\frac{2\pi nx}{3} \right|_0^2
	-\left.\frac{3}{2\pi n}(x-3)\cos\frac{2\pi nx}{3} \right|_2^3 
	+\left.\frac{9}{(2\pi n)^2}(x-3)\sin\frac{2\pi nx}{3} \right|_2^3\right) =	\\[12pt]
	&=& \displaystyle-\frac{1}{\pi n}-\frac{3}{2\pi^2 n^2}\sin\frac{4\pi n}{3}.
\end{array}
$$
Применив теорему Дирихле о поточечной сходимости ряда Фурье, видим, что построенный ряд Фурье сходится 
к периодическому (с периодом $T=3$) продолжению исходной функции при всех $x\ne 3n$, и 
$S(3n)=-1/2$ при $n=0,\pm1,\pm2,\ldots$, где $S(x)$ --- сумма ряда Ферье. 
График функции $S(x)$ имеет следующий вид
\begin{center}
	\begin{tikzpicture}
	\begin{axis}[xmin=-6, xmax=6, ymin=-1, ymax=0.5,
	width=0.8\textwidth,
	height=0.4\textwidth,
	axis x line=middle,
	axis y line=middle, 
	every axis x label/.style={at={(current axis.right of origin)},anchor=west},
	every inner x axis line/.append style={|-latex'},
	every inner y axis line/.append style={|-latex'},
	minor tick num=1,			
	axis equal=true,
	xlabel=$x$, 
	ylabel=$S(x)$,          
	samples=100,
	clip=true,
	]
	\addplot[color=black, line width=1.5pt,domain=-6:-4] {-1};
	\addplot[color=black, line width=1.5pt,domain=-4:-3]{\x+3};
	\addplot[color=black, line width=1.5pt,domain=-3:-1] {-1};
	\addplot[color=black, line width=1.5pt,domain=-1:0]{\x};
	\addplot[color=black, line width=1.5pt,domain=0:2] {-1};
	\addplot[color=black, line width=1.5pt,domain=2:3]{\x-3};
	\addplot[color=black, line width=1.5pt,domain=3:5] {-1};
	\addplot[color=black, line width=1.5pt,domain=5:6]{\x-6};
	\addplot[thick,dashed] coordinates {(-4,0) (-4,-1)};
	\addplot[thick,dashed] coordinates {(-1,0) (-1,-1)};
	\addplot[thick,dashed] coordinates {(2,0) (2,-1)};
	\addplot[thick,dashed] coordinates {(5,0) (5,-1)};
	\addplot[
	mark=*,
	mark options={fill=black, draw=black},
	only marks,
	] coordinates {(-6, -0.5) (-3, -0.5) (0, -0.5) (3, -0.5) (6, -0.5)};
	\end{axis}
	\end{tikzpicture}
\end{center}
\noindent
\textbf{Ответ:}
\[
\begin{split}
&f(x)=-\frac{5}{6}+\sum_{n=1}^\infty\left[ \frac{3}{2\pi^2 n^2}\left(1-\cos\frac{4\pi n}{3} \right)\cos\frac{2\pi nx}{3}-\left(\frac{1}{\pi n}+\frac{3}{2\pi^2 n^2}\sin\frac{4\pi n}{3} \right)\sin\frac{2\pi nx}{3}\right], x\ne 3n; \\
&S(3n)=-\frac{1}{2}, \text{ при } n\in\mathbb{Z}.
\end{split}
\]




% ---------------------------- Problem 2----------------------------------
\subsubsection*{\center Задача № 2.}
{\bf Условие.~}
Для заданной графически функции $y(x)$ построить ряд Фурье в комплексной форме, изобразить график суммы построенного ряда

%График
\begin{center}
	\begin{tikzpicture}
	\begin{axis}[xmin=-1,	xmax=3.5, 	ymin=-1,	ymax=0.5,
	width=0.5\textwidth,
	height=0.4\textwidth,
	axis x line=middle,
	axis y line=middle, 
	every axis x label/.style={at={(current axis.right of origin)},anchor=west},
	every inner x axis line/.append style={|-latex'},
	every inner y axis line/.append style={|-latex'},
	minor tick num=1,			
	axis equal=true,
	xlabel=$x$, 
	ylabel=$y$,          
	samples=100,
	clip=true,
	]
	\addplot[color=black, line width=1.5pt,domain=0:2] {-1};
	\addplot[color=black, line width=1.5pt,domain=2:3]{\x-3};
	\addplot[thick,dashed] coordinates {(2,0) (2,-1)};
	\addplot[
	mark=*,
	mark options={fill=black, draw=black},
	only marks,
	] coordinates {(2, -1)};
	\end{axis}
	\end{tikzpicture}
\end{center}

\noindent
\textbf{Решение.}\\

\noindent
Ряд Фурье в комплексной форме имеет следующий вид
\[
f(x) = \sum_{n=-\infty}^\infty c_n e^{i\omega nx},\quad c_n=\frac{1}{T}\int\limits_a^b f(x) e^{-i\omega nx}dx,\,\omega=\frac{2\pi}{T}.
\]
В нашем примере $ a=0,b=3,T=3,\omega=2\pi/3$, 
найдем коэффицинеты $c_n,\,n=0,\pm1,\pm2,\ldots$
где $\omega=2\pi/T,\,T=3.$
$$
\begin{array}{rcl}
c_0 &=&\displaystyle\frac{1}{3} \int\limits_0^3 f(x)dx=\frac{a_0}{2}=-\frac{5}{6},\\[12pt]
c_n &=&\displaystyle\frac{1}{3}\left(
-\int\limits_0^2
e^{-i\omega nx}dx+ \int\limits_2^3
(x-3) e^{-i\omega nx}dx \right) ={}\\[12pt]
&=&\displaystyle\frac{1}{3}\left(
-\left.\frac{3i}{2\pi n} e^{-i\omega nx} \right|_0^2
+\left.\frac{3i}{2\pi n}\left[(x-3) e^{-i\omega nx} - \frac{3i}{2\pi n} e^{-i\omega nx} \right]\right|_2^3\right) = \\[12pt]
&=&\displaystyle\frac{i}{2\pi n}+\frac{3}{4\pi^2 n^2}\left(1- e^{-2\omega ni} \right)=\frac{3}{4\pi^2 n^2}\left(1-\cos\frac{4\pi n}{3} \right)+\frac{i}{2\pi n}\left(1+\frac{3}{2\pi n}\sin\frac{4\pi n}{3} \right).
\end{array}
$$
\noindent
Применив теорему Дирихле о поточечной сходимости ряда Фурье, видим, что построенный ряд Фурье сходится 
к периодическому (с периодом $T=3$) продолжению исходной функции при всех $x\ne 3n$, и $S(3n)=-1/2$ при 
$n=0,\pm1,\pm2,\ldots$, где $S(x)$ --- сумма ряда Фурье. График функции $S(x)$ имеет вид
\begin{center}
	\begin{tikzpicture}
	\begin{axis}[xmin=-6, xmax=6, ymin=-1, ymax=0.5,
	width=0.8\textwidth,
	height=0.4\textwidth,
	axis x line=middle,
	axis y line=middle, 
	every axis x label/.style={at={(current axis.right of origin)},anchor=west},
	every inner x axis line/.append style={|-latex'},
	every inner y axis line/.append style={|-latex'},
	minor tick num=1,			
	axis equal=true,
	xlabel=$x$, 
	ylabel=$S(x)$,          
	samples=100,
	clip=true,
	]
	\addplot[color=black, line width=1.5pt,domain=-6:-4] {-1};
	\addplot[color=black, line width=1.5pt,domain=-4:-3]{\x+3};
	\addplot[color=black, line width=1.5pt,domain=-3:-1] {-1};
	\addplot[color=black, line width=1.5pt,domain=-1:0]{\x};
	\addplot[color=black, line width=1.5pt,domain=0:2] {-1};
	\addplot[color=black, line width=1.5pt,domain=2:3]{\x-3};
	\addplot[color=black, line width=1.5pt,domain=3:5] {-1};
	\addplot[color=black, line width=1.5pt,domain=5:6]{\x-6};
	\addplot[thick,dashed] coordinates {(-4,0) (-4,-1)};
	\addplot[thick,dashed] coordinates {(-1,0) (-1,-1)};
	\addplot[thick,dashed] coordinates {(2,0) (2,-1)};
	\addplot[thick,dashed] coordinates {(5,0) (5,-1)};
	\addplot[
	mark=*,
	mark options={fill=black, draw=black},
	only marks,
	] coordinates {(-6, -0.5) (-3, -0.5) (0, -0.5) (3, -0.5) (6, -0.5)};
	\end{axis}
	\end{tikzpicture}
\end{center}

\noindent
\textbf{Ответ:}
\[
\begin{split}
&f(x)=\sum_{n=-\infty}^\infty\left[ \frac{3}{4\pi^2 n^2}\left(1-\cos{\frac{4\pi n}{3}} \right)+\frac{i}{2\pi n}\left(1+\frac{3}{2\pi n}\sin{\frac{4\pi n}{3}} \right)\right] e^{\tfrac{i2\pi nx}{3}},~ x\ne 3n; \\
&S(3n)=-\frac{1}{2},\quad\text{при}~n\in\mathbb{Z}.
\end{split}
\]


% ---------------------------- Problem 3----------------------------------
\subsubsection*{\center Задача № 3.}
{\bf Условие.~}\\
Найти резольвенту для интегрального уравнения Вольтерры со следующим ядром 
$$K(x,t)=x.$$

\noindent
{\bf Решение.~}\\
\noindent
Запишем интегральное уравнение Вольтерры
$$
y(x)=x^3+\int\limits_0^x x y(t)dt.
$$
Из рекурентных соотношений получаем
$$
\begin{array}{rcl}
K_1(x,t)&=&\displaystyle x, \\[12pt]
K_2(x,t)&=&\displaystyle\int\limits_t^x K(x,s)K_1(s,t)ds = \int\limits_t^x x s ds = x\cdot\frac{x^2-t^2}{2},\\[12pt]
K_3(x,t)&=&\displaystyle\int\limits_t^x K(x,s)K_2(s,t)ds = \int\limits_t^x x s \frac{s^2-t^2}{2} ds = \frac{x}{2}\left(\frac{x^2-t^2}{2}\right)^2.\\[12pt]
K_j(x,t)&=&\displaystyle\frac{x}{(j-1)!}\left(\frac{x^2-t^2}{2}\right)^{j-1}\!\!\!\!\!\!\!\!,\quad j=\mathbb{N}.
\end{array}
$$
Подставляя это выражение для итерированных ядер, найдем резольвенту
$$ 
R(x,t,\lambda)=x\sum_{j=1}^\infty \frac{\lambda^{j-1}}{(j-1)!}\left(\frac{x^2-t^2}{2}\right)^{j-1}\!\!\!\!\!\!\!\!,
\quad j=1,2,\ldots
$$



%=================================================================================================================================
\subsection{Приложения дифференциального исчисления.}
% ---------------------------- Problem 1----------------------------------
\subsubsection*{\center Задача № 1.}
{\bf Условие.~}
Разложить функцию $f(x)$ по формуле Тейлора 3--го порядка в точке $x_0=1$ с остаточным членом в форме Пеано, если
$$
f(x) = x^x.
$$
{\bf Решение.~}	
Представим функцию в виде $f(x) = e^{x\ln{x}}$. Разложим элементарные функции $e^x$ и $\ln{x}$ в точке 
Выполним замену $t=x-x_0,$ отсюда $x=t+1$. В новых переменных, получим задачу о разложении функции $f(t)=e^{(t+1)\ln{(t+1)}}$ в точке $t_0=0$. Разложим спетень экспоненты $(t+1)\ln{(t+1)}$ применяя стандартное разложение
$$
\begin{array}{rl}
\ln{(t+1)} = &t - \dfrac{t^2}{2} + \dfrac{t^3}{3} + O(t^4), \\[10pt]
(t+1)\ln{(t+1)} = &(t+1)(t - \dfrac{t^2}{2} + \dfrac{t^3}{3} + O(t^4)) =\\[10pt]
=& t - \dfrac{t^2}{2} + \dfrac{t^3}{3} + t^2 - \dfrac{t^3}{2} + \dfrac{t^4}{3} + O(t^4) = \\[10pt]
=& t + \dfrac{t^2}{2} - \dfrac{t^3}{6} + O(t^4).
\end{array}
$$
Далее, раскладывая экспоненту по Формуле Тейлора, и применяя разложение для степени, получим 
$$
\begin{array}{rl}
e^t = & 1 + t + \dfrac{t^2}{2} + \dfrac{t^3}{6} + O(t^4), \\[10pt]
e^{(t+1)\ln{(t+1)}} = & 1 + \bigl(t + \dfrac{t^2}{2} - \dfrac{t^3}{6} + O(t^4)\bigr) 
+ \dfrac{1}{2}\bigl(t + \dfrac{t^2}{2} - \dfrac{t^3}{6} + O(t^4)\bigr)^2 + \\[10pt]
+& \dfrac{1}{6}\bigl(t + \dfrac{t^2}{2} - \dfrac{t^3}{6} + O(t^4)\bigr)^3 =
1 + t + \dfrac{t^2}{2} - \dfrac{t^3}{6} + \dfrac{t}{2}+\dfrac{t^3}{2} + \dfrac{t^3}{6}+O(t^4) = \\[10pt]
=& 1 + t + t^2 + \dfrac{t^3}{2} + O(t^4).
\end{array}
$$
Возвращаясь к исходной переменной $x$, получим разложение
$$
x^x = 1 + (x-1) + (x-1)^2 + \dfrac{1}{2}(x-1)^3 + O\bigl((x-1)^4\bigr).
$$

% ---------------------------- Problem 2----------------------------------
\subsubsection*{\center Задача № 2.}
{\bf Условие.~}
Исследовать данные функции и построить их графики
$$
\begin{array}{cc}
\text{\bf(а):} & y = \dfrac{x}{x^2+1}, \\[10pt]
\text{\bf(б):} & y = \sqrt[3]{x}+\sqrt[3]{12-x}, \\[10pt]
\text{\bf(в):} & y = 2\cos{x} + \dfrac{1}{2}\cos(2x), \\[10pt]
\text{\bf(г):} & y = x + 4\arcctg{x}, \\[10pt]
\text{\bf(д):} & y = (x^2+1)e^{-x^2/2}.
\end{array}
$$
{\bf Решение.~}\\
\text{\bf(а):}
\begin{center}
	\begin{tikzpicture}
	\def\func{\x/(pow(\x,2))+1)} 
	
	\begin{axis}[
	xmin=-8.75,
	xmax=8.75, 
	ymin=-2.0,
	ymax=4.5,
	width=\textwidth,
	height=0.75\textwidth,
	axis x line=middle,
	axis y line=middle, 
	every axis x label/.style={at={(current axis.right of origin)},anchor=west},
	every inner x axis line/.append style={|-latex'},
	every inner y axis line/.append style={|-latex'},
	minor tick num=1,			
	axis equal=true,
	xlabel=$x$, 
	ylabel=$y$,          
	samples=600,
	clip=true,
	]
	\addplot[color=black, line width=1.5pt,domain=-8.5:-0.1] {\func};
	\addplot[color=black, line width=1.5pt,domain=0.1:8.5]{\func};
	\addplot[color=black, dashed, domain=-8.5:8.5]{0*\x+1};
	\end{axis}
	\end{tikzpicture}
\end{center}
\text{\bf(б):}
\begin{center}
	\begin{tikzpicture}
	\def\func{pow(\x,1/3) + pow(12-\x,1/3)} 
	
	\begin{axis}[
	xmin=-.5,
	xmax=12.5, 
	ymin=0.0,
	ymax=3.5,
	width=\textwidth,
	height=0.75\textwidth,
	axis x line=middle,
	axis y line=middle, 
	every axis x label/.style={at={(current axis.right of origin)},anchor=west},
	every inner x axis line/.append style={|-latex'},
	every inner y axis line/.append style={|-latex'},
	minor tick num=1,			
	axis equal=true,
	xlabel=$x$, 
	ylabel=$y$,          
	samples=600,
	clip=true,
	]
	\addplot[color=black, line width=1.5pt,domain=0:12] {\func};
	\addplot[
	mark=*,
	mark options={fill=black, draw=black},
	only marks,
	] coordinates {(0, 2.28943) (12, 2.28943)};
	\end{axis}
	\end{tikzpicture}
\end{center}
\text{\bf(в):}
\begin{center}
	\begin{tikzpicture}
	\def\func{2*cos(deg(\x)) + 0.5*cos(2*deg(\x))} 
	
	\begin{axis}[
	xmin=-6.5,
	xmax=6.5, 
	ymin=-1.5,
	ymax=2.5,
	width=\textwidth,
	height=0.75\textwidth,
	axis x line=middle,
	axis y line=middle, 
	every axis x label/.style={at={(current axis.right of origin)},anchor=west},
	every inner x axis line/.append style={|-latex'},
	every inner y axis line/.append style={|-latex'},
	minor tick num=1,			
	axis equal=true,
	xlabel=$x$, 
	ylabel=$y$,          
	samples=600,
	clip=true,
	]
	\addplot[color=black, line width=1.5pt,domain=-6:6] {\func};
	\end{axis}
	\end{tikzpicture}
\end{center}
\text{\bf(г):}
\begin{center}
	\begin{tikzpicture}
	\def\func{4*rad(90-atan(\x))+\x} 
	
	\begin{axis}[
	xmin=-6.5,
	xmax=6.5, 
	ymin=-2.0,
	ymax=12.5,
	width=\textwidth,
	height=0.75\textwidth,
	axis x line=middle,
	axis y line=middle, 
	every axis x label/.style={at={(current axis.right of origin)},anchor=west},
	every inner x axis line/.append style={|-latex'},
	every inner y axis line/.append style={|-latex'},
	minor tick num=1,			
	axis equal=true,
	xlabel=$x$, 
	ylabel=$y$,          
	samples=600,
	clip=true,
	]
	\addplot[color=black, line width=1.5pt,domain=-12.:-0.001] {\func};
	\addplot[color=black, line width=1.5pt,domain=0.001:10] {\func};	
	\addplot[color=black, dashed, domain=-2:10.5]{\x};
	\addplot[color=black, dashed, domain=-12.5664:2]{\x+12.5664};
	\end{axis}
	\end{tikzpicture}
\end{center}
\text{\bf(д):}
\begin{center}
	\begin{tikzpicture}
	\def\func{(pow(x,2)+1)*exp(-0.5*pow(x,2))} 
	
	\begin{axis}[
	xmin=-4.5,
	xmax=4.5, 
	ymin=-0.0,
	ymax=0.5,
	width=\textwidth,
	height=0.75\textwidth,
	axis x line=middle,
	axis y line=middle, 
	every axis x label/.style={at={(current axis.right of origin)},anchor=west},
	every inner x axis line/.append style={|-latex'},
	every inner y axis line/.append style={|-latex'},
	minor tick num=1,			
	axis equal=true,
	xlabel=$x$, 
	ylabel=$y$,          
	samples=600,
	clip=true,
	]
	\addplot[color=black, line width=1.5pt,domain=-4:4] {\func};
	\end{axis}
	\end{tikzpicture}
\end{center}

% ---------------------------- Problem 3----------------------------------
\subsubsection*{\center Задача № 3.}
{\bf Условие.~}\\
Найти наибольшее расстояние точки эллипса $x=a\cos{t},\,y=b\sin{t}$ от конца
его малой полуоси $(a>b>0)$.\\
{\bf Решение.~}\\
Рассмотрим точку на эллипсе $A(x,\,y)$, при $x>0,\,y>0$. Соединим точку эллипса $A$ 
с концами малых полуосей $M_1,\,M_2$. Чертёж к задаче представлен на рис.\,\ref{fig:01}.
В силу положительности координат точки $A$, большим из двух будет отрезок $AM_2$, а расстояние
$$
\rho = \rho(A,M_2) = \sqrt{x^2 + (y+b)^2} \xrightarrow{\hphantom{aaa} } \max
$$
по условию задачи должно быть наибольшим.
Так как $x=a\cos{t}$, а $y=b\sin{t}$, то $\rho=\rho(t)$ есть функция паметра $t$ вида
$$
\rho(t) = \sqrt{a^2\cos^2{t} + (b\sin{t}+b)^2} = \sqrt{a^2\cos^2{t} + b^2(1 + \sin{t})^2}.
$$
\begin{figure}[th!]
	\center
	\begin{tikzpicture}
	
	\draw [help lines, thin] (-4, -2) grid (4, 2);
	
	\draw[ultra thick, black!60!green] (4,-2) -- (4, 0);
	\draw[ultra thick, black!60!green] (0,-2) -- (4, -2);
	
	\draw[color=blue,very thick] (0, 0) ellipse (4 and 2);
	
	\node[] at (4.5, -1)   (b) {$b$};
	\node[] at (2, -2.5)   (a) {$a>b$};
	
	\coordinate [label=above:$M_1\text{(0;\,$b$)}$] 	(M1) 	at (0,2);
	\coordinate [label=below:$M_2\text{(0;\,$-b$)}$] 	(M2) 	at (0,-2);
	\coordinate [label=right:$A\text{($x$;\,$y$)}$] 	(A) 	at (3.77124, 0.666667);
	\coordinate [label=above:$O\text{(0;\,0)}$] 	(O) 	at (0, 0);	
	
	\foreach \Point in {(0,2), (0,-2), (3.77124, 0.666667), (0, 0)}{
		\node at \Point {\textbullet};
	}
	
	\draw[ultra thick] (0,-2) -- (3.77124, 0.666667);
	\draw[] (0,2) -- (3.77124, 0.666667);
	
	\end{tikzpicture}
	\caption{Чертёж к задаче 3}
	\label{fig:01}
\end{figure}\\
Найдём экстремальные значения функции $\rho(t)$ из решения уравнения
$$
\rho'_t = 0.
$$
Вычислим производную по переменной $t$, получим
$$
\rho'_t = \dfrac{a^22\cos{t}(-\sin{t}) + b^22(1+\sin{t})\cos{t}}{2\sqrt{a^2\cos^2{t} + b^2(1 + \sin{t})^2}} = \dfrac{2\cos{t}(-a^2\sin{t} + b^2(1+\sin{t}))}{2\rho}.
$$
Приравнивая производную к нулю, получим уравнение относително переменной $t$, решая которое, находим
$$
\begin{array}{c}
2\cos{t}(-a^2\sin{t} + b^2(1+\sin{t})) = 0, \\[4pt]
-a^2\sin{t} + b^2(1+\sin{t}) = 0,			\\[4pt]
(a^2-b^2)\sin{t} = b^2,						\\[4pt]
\sin{t} = \dfrac{b^2}{a^2-b^2}.
\end{array}
$$
То есть, при $t^{\ast} = \arcsin{\dfrac{b^2}{a^2-b^2}}$ получаем экстремальные значения для расстояния $\rho_{\max}=\rho(t^{\ast})=\max\rho(A,\,M_2)$. Из равенства для $\sin{t}$ получим следствия
$$
1+\sin{t} = \dfrac{a^2}{a^2-b^2},\quad \cos^2{t} = \dfrac{a^2(a^2-2b^2)}{(a^2-b^2)^2}.
$$
Отсюда, найдём наибольшее расстояние
$$
\begin{array}{rl}
\rho_{\max}&=\rho(t^{\ast}) = \sqrt{a^2\cos^2{t^{\ast}} + b^2(1 + \sin{t^{\ast}})^2} = 
\sqrt{a^2\dfrac{a^2(a^2-2b^2)}{(a^2-b^2)^2} + b^2\dfrac{a^4}{(a^2-b^2)^2}} =  \\[4pt]
& = \dfrac{a^2}{a^2-b^2}\sqrt{a^2-b^2} = \dfrac{a^2}{\sqrt{a^2-b^2}}.
\end{array}
$$

\newpage
\addcontentsline{toc}{section}{Список литературы}
\begin{thebibliography}{99}
\bibitem{book01} Львовский С.М. Набор и вёрстка в системе \LaTeX, 2003 c.
\bibitem{book02} {\color{brightmaroon}Добавить сюда источник.}
\bibitem{book03} {\color{brightmaroon}Добавить сюда источник.}
\end{thebibliography}

\end{document}